\documentclass[11pt]{article}

\usepackage{amsmath,amssymb}
\usepackage[margin=1in]{geometry}
\usepackage{hyperref}
\usepackage[round,authoryear]{natbib}

\title{Literature Review Detailed:\\
Key Results and Intuition}
\author{Draft notes}
\date{\today}

\begin{document}
\maketitle

\noindent
This note records, paper by paper, the key results and the intuition behind
them.  It is more detailed than the positioning note
(\texttt{literature\_view}).

%=============================================================================
\section{Public finance and deadweight loss}
%=============================================================================

\subsection{The Measurement of Waste (Harberger, 1964)}

\citet{harberger_measurement_1964} develops a linearized framework for measuring
aggregate efficiency losses from taxes and related price distortions.  The paper
emphasizes six results.

\begin{itemize}
    \item \textbf{Measuring deadweight loss ($\Delta W$) as a policy tool}
    \begin{itemize}
        \item \textbf{Key result:} Deadweight loss can be calculated with a
        linearized aggregate framework that measures efficiency changes across
        distorted markets:
        \[
        \Delta W = \frac{1}{2} \sum_{i} \sum_{j} R_{ij} T_i T_j.
        \]
        \item \textbf{Intuition:} Rather than only asking whether an allocation
        is Pareto optimal, policymakers need a quantitative measure of
        efficiency loss to weigh tax reforms, tariffs, and other real-world
        trade-offs.
    \end{itemize}

    \item \textbf{Neutrality of uniform value-added taxes}
    \begin{itemize}
        \item \textbf{Key result:} An equal percentage tax on value added across
        all productive activities yields zero net efficiency loss:
        \[
        \sum_{j} v_j R_{ij} = 0
        \implies
        \Delta W = 0
        \quad\text{when } t_i = t \text{ for all } i.
        \]
        \item \textbf{Intuition:} Uniform value-added taxation leaves relative
        price ratios across activities unchanged, so it does not induce
        distortionary factor substitution.
    \end{itemize}

    \item \textbf{Equivalence of product taxes and value-added taxes}
    \begin{itemize}
        \item \textbf{Key result:} Product taxes can be translated into
        equivalent taxes on value added when intermediate inputs enter output
        in fixed proportions.
        \item \textbf{Intuition:} Taxing final output creates the same
        behavioral incentives and equilibrium as taxing all factor inputs
        (labor, capital, materials) at the same effective rate.
    \end{itemize}

    \item \textbf{Cross-tax distortions and welfare-improving tax increases}
    \begin{itemize}
        \item \textbf{Key result:} Introducing or raising a tax on a good can
        raise total social welfare if that good is a substitute for another
        already heavily taxed good:
        \[
        \frac{\partial (\Delta W)}{\partial T_k} > 0
        \quad\text{when } T_j > 0 \text{ and } R_{kj} > 0.
        \]
        \item \textbf{Intuition:} If Good~$A$ is heavily taxed, consumers
        substitute toward Good~$B$.  Taxing Good~$B$ shifts demand back toward
        Good~$A$, reducing the preexisting distortion in Sector~$A$ by enough
        to offset the new distortion in Sector~$B$.
    \end{itemize}

    \item \textbf{Distortion of the labor--leisure choice}
    \begin{itemize}
        \item \textbf{Key result:} A uniform tax on all market labor is
        non-neutral unless leisure (non-market activity) can also be taxed:
        \[
        \sum_{j=1}^{n+1} M_{ji} = 0.
        \]
        \item \textbf{Intuition:} Leisure is hard to tax.  Taxing all market
        labor therefore induces substitution into untaxed leisure and creates
        an efficiency loss in labor supply.
    \end{itemize}

    \item \textbf{Capital--labor cross-factor interaction effects}
    \begin{itemize}
        \item \textbf{Key result:} When taxes distort both capital ($B_i$) and
        labor ($E_i$), cross-factor interaction terms
        ($H_{ij}=\partial K_i/\partial E_j$) must be included to measure total
        efficiency loss:
        \[
        \Delta W
        = \frac{1}{2}\sum_{i}\sum_{j} M_{ij} E_i E_j
        + \frac{1}{2}\sum_{i}\sum_{j} G_{ij} B_i B_j
        + \sum_{i}\sum_{j} H_{ij} B_i E_j.
        \]
        \item \textbf{Intuition:} A tax on capital in one industry reallocates
        both capital and labor across sectors.  Ignoring how labor responds to
        capital taxation (and vice versa) misstates total efficiency loss.
    \end{itemize}
\end{itemize}

\subsection{Salience and Taxation: Theory and Evidence (Chetty, Looney, \& Kroft, 2009)}

\citet{chetty_salience_2009} show that consumer responses to taxation depend on
whether the tax is salient at the point of choice.  The paper combines a field
experiment with evidence from alcohol taxation to argue that standard tax
theory misses an important channel when buyers do not fully perceive
non-salient taxes.

\begin{itemize}
    \item \textbf{Underreaction to non-salient taxes}
    \begin{itemize}
        \item \textbf{Key result:} Consumers consistently underreact to taxes
        that are not explicitly included in posted price tags.
        \item \textbf{Intuition:} Standard models assume buyers continuously
        compute the full final price.  In practice, doing so requires active
        mental effort, so many consumers fail to fully incorporate taxes that
        are added only at checkout.
    \end{itemize}

    \item \textbf{The salience parameter $\theta$}
    \begin{itemize}
        \item \textbf{Key result:} The paper summarizes attention with a
        salience parameter $\theta \in [0,1]$, so consumers behave as if they
        face the perceived price $\tilde{p}=p+\theta t$ rather than the true
        price $p+t$.
        \item \textbf{Intuition:} When $\theta=1$, the model collapses to the
        standard full-optimization benchmark and consumers treat a tax increase
        like a base-price increase.  When $\theta=0$, consumers ignore the tax
        completely while shopping.  The empirically relevant case is
        $0<\theta<1$, where consumers notice the tax but underreact to it.
    \end{itemize}

    \item \textbf{Grocery store field experiment}
    \begin{itemize}
        \item \textbf{Key result:} Posting tax-inclusive shelf tags reduced
        demand by about 8 percent.
        \item \textbf{Intuition:} Consumers rely heavily on the visible shelf
        price as a decision shortcut.  Making the tax visible at the moment of
        choice changes behavior because it converts a hidden cost into an
        immediately attended price signal.
    \end{itemize}

    \item \textbf{Excise tax versus sales tax impact}
    \begin{itemize}
        \item \textbf{Key result:} State-level increases in alcohol excise
        taxes, which are embedded in posted prices, reduced beer consumption by
        an order of magnitude more than equivalent increases in general sales
        taxes, which are added at checkout.
        \item \textbf{Intuition:} Taxes that are already incorporated in shelf
        prices affect the price consumers actually see while shopping, whereas
        register-added taxes remain less salient until after most purchase
        decisions are made.
    \end{itemize}

    \item \textbf{Economic incidence depends on statutory form}
    \begin{itemize}
        \item \textbf{Key result:} Economic incidence depends on the statutory
        form of the tax once salience matters, so excise taxes and register-added
        sales taxes need not place the same burden on consumers and producers.
        \item \textbf{Intuition:} When a tax is non-salient, consumers do not
        reduce demand as much as they would under full attention.  Because
        quantity demanded remains artificially high, producers do not need to
        cut pre-tax prices as much to sustain sales, and inattentive consumers
        bear a larger share of the burden.  This overturns the standard view
        that statutory form is irrelevant for incidence.
    \end{itemize}

    \item \textbf{Incidence formula adjustment under inattention}
    \begin{itemize}
        \item \textbf{Key result:} In the standard model, the producer burden
        depends on relative elasticities, with
        $\text{Burden on producers} \propto \epsilon_D/(\epsilon_S+\epsilon_D)$.
        In the behavioral model, the demand response to a tax is attenuated to
        $\theta \epsilon_D$.
        \item \textbf{Intuition:} Because consumers react only to the salient
        fraction of the tax, demand falls less after a tax increase than a
        standard model predicts.  Producers therefore do not need to lower
        pre-tax prices as much, so more of the tax is shifted onto consumers.
        In the limiting case $\theta=0$, demand does not respond to the tax at
        all and producers can pass through the full tax without absorbing any
        of the burden through lower prices.
    \end{itemize}

    \item \textbf{Decision miscalculation and overconsumption}
    \begin{itemize}
        \item \textbf{Key result:} Consumers choose based on a perceived price
        $\tilde{p}=p+\theta t$ with $\theta<1$, rather than the true tax-inclusive
        price $p+t$, so non-salient taxes can generate excess burden even if
        observed demand barely moves.
        \item \textbf{Intuition:} Consumers buy extra units whose marginal
        benefit is below their true marginal cost because the tax is only
        partially internalized at the time of choice.  Section~V of the paper
        formalizes this point: decision errors themselves generate deadweight
        loss, even when observed quantity responses are very small or zero.
    \end{itemize}

    \item \textbf{Deadweight-loss formula adjustment}
    \begin{itemize}
        \item \textbf{Key result:} Welfare loss from demand distortion shrinks
        with tax inattention, so a useful generalized Harberger benchmark is
        $\text{DWL}_{\text{behavioral}} \approx \tfrac12 pQ\epsilon_D \theta t^2$.
        But total consumer welfare still falls because inattentive consumers
        make costly decision errors even when observed demand changes very
        little ($\Delta Q \approx 0$).
        \item \textbf{Intuition:} Standard tax analysis measures harm mainly by
        how many purchases are canceled.  When taxes are hidden, a small
        $\Delta Q$ does not mean the tax is harmless; it means consumers are
        being induced to buy goods whose true cost exceeds their value.
        Welfare losses therefore arise through optimization mistakes and
        unexpected budget squeezes, not only through the usual demand-drop
        channel.
    \end{itemize}

    \item \textbf{Unexpected budget squeezes and suboptimal cutbacks}
    \begin{itemize}
        \item \textbf{Key result:} Non-salient taxes can worsen welfare by
        causing consumers to violate their intended budget allocation while
        browsing and then forcing unplanned adjustments at checkout.
        \item \textbf{Intuition:} Because the omitted tax is not reflected in
        the posted price, consumers do not initially choose bundles that satisfy
        their true budget constraint $Z=(p+t)x+y$.  Once the tax is charged,
        they must cut spending on other goods unexpectedly, ending up with a
        worse overall consumption bundle than they would have chosen under full
        information.
    \end{itemize}
\end{itemize}

\bibliographystyle{plainnat}
\bibliography{references}

\end{document}
